\par\Needspace{4\baselineskip}
\begin{center}
\hypertarget{native-036-article018}{}\pdfbookmark[2]{18.}{article-018}18.
\end{center}
In der PARABEL wären die exzentrische Anomalie, die mittlere Anomalie und die mittlere Bewegung sämtlich $=0$. Diese Begriffe eignen sich hier also nicht dazu, die Bewegung zur Zeit in Beziehung zu setzen. Bei der Parabel brauchen wir jedoch zur Integration von $r r\,\mathrm{d}v$ überhaupt keinen Hilfswinkel, denn es gilt
\[
\begin{aligned}
r r\,\mathrm{d}v&=\frac{p p\,\mathrm{d}v}{4\cos\frac{1}{2}v^4}=\frac{p p\,\mathrm{d}\operatorname{tang}\frac{1}{2}v}{2\cos\frac{1}{2}v^2}\\
&=\frac{1}{2}p p(1+\operatorname{tang}\frac{1}{2}v^2)\,\mathrm{d}\operatorname{tang}\frac{1}{2}v
\end{aligned}
\]
und somit $\int r r\,\mathrm{d}v=\frac{1}{2}p p(\operatorname{tang}\frac{1}{2}v+\frac{1}{3}\operatorname{tang}\frac{1}{2}v^3)+\mathrm{Const.}$ Rechnet man die Zeit vom Durchgang durch das Perihel an, so ist die Konstante $=0$. Damit erhält man
\[
\operatorname{tang}\frac{1}{2}v+\frac{1}{3}\operatorname{tang}\frac{1}{2}v^3=\frac{2t k\BareRoot{(1+\mu)}}{p^{\frac{3}{2}}},
\]
\SourceLeaf{037}Mit dieser Formel lässt sich $t$ aus $v$ und $v$ aus $t$ bestimmen, sobald $p$ und $\mu$ bekannt sind. Unter den parabolischen Elementen gibt man gewöhnlich anstelle von $p$ den Radiusvektor im Perihel an, der $\frac{1}{2}p$ beträgt, und vernachlässigt die Masse $\mu$ ganz. Die Masse eines Körpers, dessen Bahn als Parabel berechnet wird, wird sich gewiss kaum je bestimmen lassen. Tatsächlich scheinen nach den besten und neuesten Beobachtungen alle Kometen eine so geringe Dichte und Masse zu haben, dass ihre Masse als unmerklich gelten und ohne Bedenken vernachlässigt werden kann.
